% Generated draft: requires mathematical and source review.
\documentclass[11pt]{article}
\usepackage[margin=25mm]{geometry}
\usepackage{fontspec}
\setmainfont{Latin Modern Roman}
\usepackage{amsmath,amssymb,mathtools,unicode-math}
\setmathfont{Latin Modern Math}
\usepackage{enumitem,needspace,xurl,hyperref}
\setlength{\parindent}{0pt}
\setlength{\parskip}{5pt}
\setlength{\emergencystretch}{3em}
\newcommand{\N}{\mathbb{N}}
\newcommand{\Z}{\mathbb{Z}}
\newcommand{\Q}{\mathbb{Q}}
\newcommand{\R}{\mathbb{R}}
\newcommand{\C}{\mathbb{C}}
\newcommand{\F}{\mathbb{F}}
\newcommand{\A}{\mathbb{A}}
\newcommand{\PP}{\mathbb{P}}
\newcommand{\E}{\mathbb{E}}
\newcommand{\eps}{\varepsilon}
\newcommand{\dd}{\,\mathrm{d}}
\begin{document}
\section*{UnsolvedMath ID: 20001854}
\textbf{Source problem code:} catalogue entry
\subsection{Definitions and mathematical statement}
Sitharam's source asks for a Farber-style Morse-theoretic calculation of chambers and homology, including Betti numbers, for two polygons sharing a chain and more generally graphs of tree-width two. The aid describes the omitted source figure as a theta graph whose cycles share a four-edge path. The original text does not specify ambient dimension, whether crossings are allowed, the symmetry quotient, coefficient ring, or the meaning of chambers.

For the aid's planar special case, let \(G=(V,E)\) be a labelled graph with positive edge lengths \(\ell_e\), and consider placements \(p:V\to\R^2\) satisfying \(|p(v)-p(w)|=\ell_{vw}\), modulo orientation-preserving Euclidean motions \(\operatorname{SE}(2)\). Crossings and coincidences of nonadjacent vertices are allowed; no quotient by reflection is taken. Pin a chosen positive-length edge to \((0,0),(c,0)\) to obtain the normalized realization space \(\mathcal C(G,\ell)\).

If \(G_1,G_2\) share exactly a path \(P\), normalize all three spaces using the same first edge of \(P\). Restriction to the path then gives the exact fiber-product description
\[
 \mathcal C(G_1\cup_PG_2,\ell)
 \cong\mathcal C(G_1,\ell)\times_{\mathcal C(P,\ell)}
                           \mathcal C(G_2,\ell).
\]
When \(P\) is a single edge, its normalized space is a point and this is an ordinary product. For a flexible multi-edge path the fiber product cannot in general be replaced by a product.

In particular, suppose two quadrilateral cycles share one edge of length \(c>0\), and the complementary three-edge arm of cycle \(j\) has lengths \((a_j,b_j,d_j)\), all positive. If, for \(j=1,2\),
\[
 |a_j-b_j|<|c-d_j|<a_j+b_j<c+d_j,
\]
then the aid reports that each normalized quadrilateral space is a circle, and therefore
\[
 \mathcal C(G,\ell)\cong S^1\times S^1=T^2.
\]
With integral homology, its Betti numbers are \(b_0=1\), \(b_1=2\), \(b_2=1\), and zero in higher degrees. The strict inequalities specify an open length family; this is a special calculation rather than a theorem for the four-edge shared path in the figure.

For comparison, let \(G\) be a maximal \(2\)-tree with \(n\geq2\) labelled vertices, built from an edge by successively adding a vertex adjacent to both endpoints of an existing edge. If every construction triangle satisfies strict triangle inequalities, each step has two distinct planar choices. In the same crossing-allowed, reflection-distinguishing normalization, the aid gives exactly
\[
 |\mathcal C(G,\ell)|=2^{n-2},\qquad
 H_0(\mathcal C(G,\ell);\Z)\cong\Z^{2^{n-2}},\quad
 H_q(\mathcal C(G,\ell);\Z)=0\ (q>0).
\]
This finite rigid case and the torus family do not supply the requested general chamber decomposition or a general Morse-theoretic homology formula for all partial \(2\)-trees.
\subsection{Short English statement}
The aid models planar crossing-allowed linkages as fiber products over their shared chain. A rigid shared edge makes this a product, giving an explicit torus family with Betti numbers \((1,2,1)\). Strictly realizable maximal \(2\)-trees instead have finitely many normalized realizations; the general chamber and homology problem remains.
\subsection{Sources}
\begin{itemize}
\item Local UnsolvedMath registry, permanent ID \(20001854\), source code \texttt{catalogue entry}; canonical collection \texttt{aim-geometry-notes.json}, notes[191].
\item AIM workshop ``Configuration spaces of linkages,'' problem \(5\), attributed in the source to Sitharam: \url{https://aimath.org/pastworkshops/linkagesproblems.pdf}. The source says ``Apply Morse Theory \`a la Farber'' for chambers and homologies. The aid describes the source figure's four-edge shared chain; the figure is not included in this record and was not opened in this drafting session.
\item Supplied research attempt \(1\), result \texttt{partial\_result}, run status \texttt{partially\_solved}, rigor label ``proved,'' novelty confidence ``low.'' Registry classification \texttt{PARTIAL-PROGRESS}, updated 14 August 2026. Original remarks and literature fields are empty; no complete authoritative references, proposer-field value or proposed year are supplied despite the Sitharam attribution.
\end{itemize}
\end{document}
